\documentclass[a4,11pt]{aleph-notas}
% Se puede ver la documentación aquí:
% https://github.com/alephsub0/LaTeX_aleph-notas

% -- Paquetes adicionales
\usepackage{enumitem}
\usepackage{url}
\usepackage{array}
\usepackage{booktabs}
\usepackage{longtable}
\usepackage{aleph-comandos}
\hypersetup{
    urlcolor=blue,
    linkcolor=blue,
}


% -- Datos
\institucion{Facultad de Ciencias Exactas, Naturales y Ambientales}
\carrera{Catálogo STEM}
\asignatura{Matemática III}
\tema{Plan tema 06: Derivadas de orden superior}
\autor{Andrés Merino}
\fecha{Periodo 2026-1}

\logouno[0.16\textwidth]{Logos/logoPUCE_04_ac}
\definecolor{colortext}{HTML}{1957d1}
\definecolor{colordef}{HTML}{1957d1}
\fuente{montserrat}


\begin{document}

\encabezado

%%%%%%%%%%%%%%%%%%%%%%%%%%%%%%%%%%%%%%%%
\section*{Resultado de Aprendizaje}
%%%%%%%%%%%%%%%%%%%%%%%%%%%%%%%%%%%%%%%%

%%%%%%%%%%%%%%%%%%%%%%%%%%%%%%%%%%%%%%%%
\subsection*{RdA de la asignatura:}
\begin{itemize}[leftmargin=*]
    \item \textbf{RdA 1:} Identificar conceptos, teoremas y propiedades de funciones vectoriales, derivadas parciales e integral vectorial.
    \item \textbf{RdA 2:} Calcular derivadas parciales e integrales con asistentes computacionales.
\end{itemize}

%%%%%%%%%%%%%%%%%%%%%%%%%%%%%%%%%%%%%%%%
\subsection*{RdA del tema:}
\begin{itemize}[leftmargin=*]
    \item Calcular derivadas parciales de segundo orden y derivadas mixtas de un campo escalar.
    \item Construir la matriz hessiana de un campo escalar en un punto.
    \item Aplicar el teorema de simetría para determinar cuándo las derivadas mixtas coinciden.
    \item Obtener la aproximación lineal y cuadrática de un campo escalar mediante las fórmulas de Taylor.
\end{itemize}

%%%%%%%%%%%%%%%%%%%%%%%%%%%%%%%%%%%%%%%%
\section*{Introducción}
%%%%%%%%%%%%%%%%%%%%%%%%%%%%%%%%%%%%%%%%

\paragraph{Pregunta inicial:}
Si conoces cómo varía la temperatura de una habitación al desplazarte en distintas direcciones, ¿podrías predecir con más precisión cómo cambiaría si te mueves en dos direcciones a la vez? ¿Qué información adicional necesitarías?

%%%%%%%%%%%%%%%%%%%%%%%%%%%%%%%%%%%%%%%%
\section*{Desarrollo}
%%%%%%%%%%%%%%%%%%%%%%%%%%%%%%%%%%%%%%%%

%%%%%%%%%%%%%%%%%%%%%%%%%%%%%%%%%%%%%%%%
\subsection*{Actividad 1: Derivadas mixtas y matriz hessiana}
Se extiende la noción de derivada parcial a órdenes superiores, se introduce la simetría de las derivadas mixtas y se construye la matriz hessiana.

\paragraph{¿Cómo lo haremos?}
\begin{itemize}[leftmargin=*]
    \item \textbf{Motivación:} partiendo de las derivadas parciales conocidas, plantear la pregunta de qué ocurre al derivar nuevamente respecto a otra variable; mostrar con un ejemplo concreto que el orden puede o no importar.
    \item \textbf{Clase magistral:} se define la derivada parcial mixta $D_{j,i}f$; se enuncia el teorema de simetría (coincidencia de derivadas mixtas bajo continuidad); se construye la matriz hessiana como representante de la segunda derivada. Se usa el resumen \href{https://andres-merino.github.io/MatematicaIII-05-N0051/01Resumen/Resumen06.pdf}{Resumen06.pdf}.
    \item \textbf{Implementación computacional:} calcular matrices hessianas de varios campos escalares con un asistente computacional y verificar la simetría.
    \item \textbf{Resolución de ejercicios:} calcular derivadas de orden superior y construir matrices hessianas, usando \href{https://andres-merino.github.io/MatematicaIII-05-N0051/04Ejercicios/Ejercicios06.pdf}{Ejercicios06.pdf}.
\end{itemize}

\paragraph{Verificación de aprendizaje:}
\begin{itemize}[leftmargin=*]
    \item Para $f(x,y)=x^3y^2$, ¿cuánto vale $D_{1,2}f$ y $D_{2,1}f$ en $(1,1)$? ¿Coinciden?
    \item ¿Qué condición garantiza que $D_{1,2}f = D_{2,1}f$?
    \item ¿Cuál es la matriz hessiana de $f(x,y)=e^x\sen(y)$ en el origen?
\end{itemize}

%%%%%%%%%%%%%%%%%%%%%%%%%%%%%%%%%%%%%%%%
\subsection*{Actividad 2: Fórmulas de Taylor}
Se presentan las aproximaciones lineal y cuadrática de un campo escalar a través de las fórmulas de Taylor de primer y segundo orden.

\paragraph{¿Cómo lo haremos?}
\begin{itemize}[leftmargin=*]
    \item \textbf{Clase magistral:} se enuncian las fórmulas de Taylor de primer y segundo orden; se interpreta la aproximación lineal (plano tangente) y la aproximación cuadrática como mejoras sucesivas; se muestra el papel de la hessiana en la aproximación cuadrática. Se usa el resumen \href{https://andres-merino.github.io/MatematicaIII-05-N0051/01Resumen/Resumen06.pdf}{Resumen06.pdf}.
    \item \textbf{Resolución de ejercicios:} obtener aproximaciones lineales y cuadráticas de campos escalares dados, usando \href{https://andres-merino.github.io/MatematicaIII-05-N0051/04Ejercicios/Ejercicios06.pdf}{Ejercicios06.pdf}.
\end{itemize}

\paragraph{Verificación de aprendizaje:}
\begin{itemize}[leftmargin=*]
    \item ¿Cuál es la aproximación lineal de $f(x,y)=\cos(x)\cos(y)$ cerca del origen?
    \item ¿Cuál es la aproximación cuadrática de $f(x,y)=e^{x+y}$ cerca del origen?
    \item ¿En qué se diferencia la aproximación cuadrática de la lineal cuando nos alejamos del punto de expansión?
\end{itemize}

%%%%%%%%%%%%%%%%%%%%%%%%%%%%%%%%%%%%%%%%
\section*{Cierre}
%%%%%%%%%%%%%%%%%%%%%%%%%%%%%%%%%%%%%%%%

\paragraph{Tarea:}
Resolver del libro \href{https://catalogobiblioteca.puce.edu.ec/cgi-bin/koha/opac-detail.pl?biblionumber=83405&query_desc=su%3A%22An%C3%A1lisis%20vectorial%22}{Cálculo Vectorial de Colley}: sección 2.4, ejercicios 1--21 (impares) y 22; sección 2.3, ejercicios 43, 44, 45.

\paragraph{Pregunta de investigación:}
\begin{enumerate}[leftmargin=*]
    \item ¿Cómo se generaliza la fórmula de Taylor a orden $n$ para funciones de varias variables?
    \item ¿En qué aplicaciones de la física o la ingeniería se usan las aproximaciones de Taylor de campos escalares?
\end{enumerate}

\paragraph{Para el próximo tema:}
Ver los videos:
\begin{itemize}[leftmargin=*]
    \item \href{https://youtu.be/5mZeT0-fzS0?si=nTnQPSBdLX6ptFpP}{Derivadas parciales de campos vectoriales}.
    \item \href{https://youtu.be/NO3AqAaAE6o?si=bZUZk__l-ebqPfhN}{Multivariable chain rule}.
\end{itemize}


\end{document}
